% Working title: Structure-Preserving Equation Discovery for Power-System
% Transients from Partial Observations
\section{Introduction}

Following a severe disturbance, synchronous generators can swing apart and
lose synchronism~\cite{hatziargyriou2021stability}. The trajectories of rotor
angles and speeds show which machines separate, how quickly the separation
develops, and whether the relative motion settles. Detailed dynamic models
describe this evolution by representing the network, generators, and their
controls. Repeated studies over many contingencies also call for compact
dynamic equivalents that can be integrated efficiently and whose equations can
be inspected. Multi-event identification has shown how disturbances at
different locations and of different types can inform a common dynamic
equivalent~\cite{jiang2020multievent}. We extend this question from fitting an
equivalent of prescribed form to discovering an explicit equation for
post-clearing motion. Such an equation can generate the subsequent transient and
expose the forces and energy exchanges that produce it.

The main difficulty is that the observed rotor angles and speeds need not form
a closed dynamical state. Two events can reach the same angles and speeds
under the same external input and still have different rotor accelerations
because their internal electrical states differ. A model of the form
\(\dot y=f(y,u)\), where \(y\) is the observed state and \(u\) is the applied
input, assigns both events the same derivative at that instant. The
Mori--Zwanzig formalism shows how eliminating dynamic variables can instead
produce both a memory term driven by the history of the retained variables and
a contribution from unresolved initial conditions~\cite{gouasmi2017memory,pagnier2025kron}.
The modeling problem is therefore one of dynamic closure. The reduced state
must retain enough information about the omitted processes to determine the
subsequent evolution under the specified input.

Physical reduction supplies a starting point for constructing that state.
Kron reduction uses a Schur complement to eliminate algebraic network nodes
while preserving the electrical relation at the retained
nodes~\cite{dorfler2013kron}. The resulting effective couplings can inform a
generator-level model. The effect of an omitted device state with its own
evolution requires a separate closure. Parameter identification within a
specified dynamical structure can estimate coupling and damping coefficients
or adjust the internal initial condition, as in structure-preserving
calibration of port-Hamiltonian models~\cite{gunther2024calibration}. The
specified structure determines which mechanisms these parameters can
represent. If a restoring force, dissipation law, or memory process is absent
from the model family, fitting can shift its effect into the
terms that are available. This motivates selecting dynamical terms within a
physically constrained candidate set.

Sparse identification of nonlinear dynamics (SINDy) provides a way to select
equation terms and estimate their coefficients from a library of candidate
functions~\cite{brunton2016sindy}. Weak formulations replace pointwise
derivative estimates with integrals against test functions, reducing the
dependence on numerical differentiation~\cite{messenger2021weak}. In power systems,
symbolic regression has already incorporated trigonometric terms, measured
external signals, and physical constraints to refine synchronous-generator
models~\cite{saric2021symbolic}. With partial observations, however, some
candidate terms depend on states that are absent from the data. A library
constructed from angles and speeds alone can attribute a memory effect to an
instantaneous term that happens to correlate with it on the observed
trajectories. A more accurate estimate of the observed derivative leaves this
ambiguity intact. Structure selection must therefore account for both the
visible terms and the dynamics that close the equations.

Prior work offers several ways to represent missing dynamics. Deep delay
autoencoders can learn effective coordinates and sparse equations together
from time-delay embeddings of partial measurements~\cite{bakarji2023delay}.
Here the generator angles and speeds have specific physical meanings that
the reduced model must preserve. Data-driven memory models describe how past
observations influence the present~\cite{lin2021wiener}. To use such a closure
in a compact continuous-time simulator, one must choose which memory terms to
realize through auxiliary states and determine those states at the forecast
origin. Port-Hamiltonian learning provides a way to encode energy exchange,
dissipation, and subsystem interconnection in the resulting
equations~\cite{neary2023compositional}. Together these ideas motivate a
formulation in which the candidate memory terms already have a finite-dimensional
realization and interact with the retained rotor motion through specified
energy exchanges. The resulting closure can then be assessed in the physical
coordinates used to measure and predict the transient.

The central question is whether restoring forces, dissipation, and memory can
be identified jointly from several partially observed transients as one
shared, physically structured ordinary differential equation (ODE) that can
be initialized for a new event using a short record. Shared coefficients
determine how the model evolves across events, and event-specific initial
memory determines its starting state. A close fit to one observed interval
can arise from several combinations of the two. Across events, complementary
excitation can help distinguish them, provided the data excite the relevant
effects. An equation fitted to observed trajectories must also be able to
predict beyond the segment used to infer its initial state. We use the first
100~ms after clearing as the observation window for this forecast test.

We formulate the task as offline equation discovery followed by causal
forecasting. Offline, the available portions of the \texttt{train} trajectories
inform the selection of restoring, dissipative, and memory terms. The fit uses
the observed trajectories to learn a single equation across events while
accounting for each event's unobserved initial memory.
The aim is to retain supported physical terms and represent the selected
memory by auxiliary equations that evolve alongside the observed rotor
coordinates. At prediction time, the model parameters and active terms are
fixed. The short history supplies the observed state and is used to infer the
memory state of the same equation. For the benchmark forecasts, we hold the
mechanical-power deviation at its forecast-origin value. From that point,
the ODE generates the future transient without updates from later observations.

The proposed study combines three contributions. First, it formulates
structure selection in terms of shared restoring, dissipative, and memory
components whose effects must be distinguished from each event's initial
memory. Second, it couples a finite-dimensional realization with an explicit
energy balance to a short-history initializer for the same ODE. Third, it
defines tests of function recovery, dynamic response, and initialization on
known-truth systems, with cross-disturbance trajectory prediction as the main
test on higher-order simulations of the IEEE 39-bus system generated with
ANDES~\cite{cui2021andes}. The known-truth controls help distinguish inadequate
excitation from errors in the chosen dynamics or initial state. Rotor-angle
stability judgments derived from the forecast trajectories provide an
additional engineering check. Section~II defines the problem and physical
constraints, and Section~III describes discovery, realization, and
initialization. Section~IV presents the validation and power-system case study,
and Section~V discusses the scope of the results.
